% New source after checkpoint08.  Root reserved A.4.47.
\subsection{An explicit finite recognition quotient for the translation normalizer}

Retain the actual normalized configuration space $X$, its grid basepoint,
and $K_0=\pi_1(X)$ from A.4.46.  Put
$\mathcal M=\Gamma_{1,[9]}$ and let $H\cong C_3^2$ be the named
translation subgroup.  In this subsection the compact period is written
in the aligned basis (MFP.1)--(MFP.4).  Thus its matrix is
$A=JC_{\mathrm{up}}J$, the grid labels are the vectors of
$\mathbb F_3^2$, and the established finite normalizer comparison reads
$\nu=A\bmod3$.

Let
\[
 q:\mathcal M\twoheadrightarrow Q=G^\dagger
                                                               \tag{ECR.1}
\]
be the actual joint quotient (FCD.2).  All its original source maps,
including the nine ordinary forgetting coordinates, remain part of
(ECR.1).  Write $D=|Q|$ and $I=24\cdot9!$.  The construction below also
works after joining additional specified finite source coordinates:
one uses the order of that actual joint image in place of $D$.

Define increasing integer functions, for $r\geq2$, by
\[
 \begin{split}
 A_3(r)&=3^{\,r(r+1+3^r(r-1))},\\
 C(1,r)&=3^r,\\
 a&=C(l-1,r)A_3(r),\\
 C(l,r)&=a\,C\bigl(l-1,1+a(r-1)\bigr)
                    A_3\bigl(1+a(r-1)\bigr)\quad(l\geq2).
 \end{split}                                                    \tag{ECR.2}
\]
Starting with $D_0=D$, set, for $0\leq i<289$,
\[
 m_i=D_iC(8,1+8D_i),\qquad
 D_{i+1}=(m_i!)^{(m_i!)^{44}}.
                                                               \tag{ECR.3}
\]
These are deliberately coarse bounds on explicitly defined finite
quotients, rather than an unspecified limiting process.

\begin{result}{A.4.47}{Proved}{A finite characteristic configuration refinement recognizes the canonical translation normalizer}
There is a specified characteristic subgroup $V\leq K_0$ of index at
most $D_{289}$ and a specified finite quotient
\[
 \eta:\mathcal M\twoheadrightarrow E,\qquad
 p:E\twoheadrightarrow Q,\qquad p\eta=q,
                                                               \tag{ECR.4}
\]
with
\[
 \ker\eta\cap K_0=V,\qquad
 |E|\leq I\bigl(D D_{289}^{45}\bigr)^I.
                                                               \tag{ECR.5}
\]
The named subgroup $H$ embeds in $E$.  The projection of its actual
finite normalizer is exactly the ordinary normalizer image:
\[
 \begin{split}
 p\bigl(N_E(\eta H)\bigr)&=q\bigl(N_{\mathcal M}(H)\bigr),\\
 p\bigl(C_E(\eta H)\bigr)&=q\bigl(C_{\mathcal M}(H)\bigr).
 \end{split}                                                    \tag{ECR.6}
\]
In particular the two maps $d\operatorname{pr}_G$ and $j_0^\dagger$
agree on the zero-label part of the image in the first line of
(ECR.6), with the aligned target marking of (MFP.3).
This conclusion uses the exact new projection (ECR.4); it does not
assert the equality on the unrefined normalizer disproved in A.4.42.
\end{result}

\begin{proof}
\par\medskip\noindent\textbf{The source group and its finite-index split subgroup.}\quad
The ordered torus configuration evaluation map and its translation
factor in (NCF.1) identify the point-pushing subgroup with
\[
 K_0=\ker\bigl(P\Gamma_{1,9}\longrightarrow
                              \operatorname{SL}_2(\mathbb Z)\bigr).
                                                               \tag{ECR.7}
\]
This kernel is normal in the full mixed group: it is the simultaneous
kernel of the label permutation and the compact period.  The actual
successive forgetful maps on $X$ give eight free extensions, of ranks
$2,3,\ldots,9$.  They give a set of $44$ generators of $K_0$.
An index-$d$ subgroup inherits these extension maps; its free kernel
ranks are at most $1+8d$, by Schreier's formula applied to the actual
intersections with the free kernels.

Let $\Gamma(3)$ denote the kernel of reduction
$\operatorname{SL}_2(\mathbb Z)\to\operatorname{SL}_2(\mathbb F_3)$.
Then
\[
 \mathcal U=P\Gamma_{1,9}\cap\operatorname{per}^{-1}\Gamma(3)
              =K_0\rtimes s\Gamma(3),\qquad
 [\mathcal M:\mathcal U]=I,
                                                               \tag{ECR.8}
\]
where $s(A)$ is the actual linear map of the torus with period $A$.
Such a map fixes every grid point because $A\equiv I\pmod3$;
linear composition proves that $s$ is a homomorphic section.
Every label permutation is represented in a disc with period one,
and the ordinary pure forgetting map is onto the torus mapping class
group.  Hence the joint label and modulo-three period image is
$\mathfrak S_9\times\operatorname{SL}_2(\mathbb F_3)$, proving the
index in (ECR.8).  Also $s\Gamma(3)$ commutes with every named
translation.  For an elementary description of its source,
$\Gamma(3)$ injects into
$\operatorname{PSL}_2(\mathbb Z)=C_2*C_3$ and has image the kernel
onto $\operatorname{PSL}_2(\mathbb F_3)=A_4$.
Its Bass--Serre quotient has twelve edges, six vertices of one type
and four of the other, so it is free of rank three.

\par\medskip\noindent\textbf{An explicit restriction that kills positive cohomology.}\quad
All cohomology in the next three paragraphs has coefficients
$k=\mathbb F_3$.  For a free group $F$ of rank $r\geq2$, put
\[
 N=F'F^3,\qquad Q_F=F/(N'N^3),\qquad
 |Q_F|=3^{r+1+3^r(r-1)}.
                                                               \tag{ECR.9}
\]
The regular graph-cover centralizer calculation of A.4.7 gives
$Z(Q_F)\to F/N$ the trivial image.  In particular this assertion
concerns the actual quotient map in (ECR.9).
For any group $G$ in which this $F$ is normal, its conjugation action
gives a homomorphism $G\to\operatorname{Aut}(Q_F)$, since $N'N^3$
is characteristic in $F$.  Its kernel $J$ satisfies
$J\cap F\subset N$.  The $r$ generator images determine every
automorphism of $Q_F$, giving
$[G:J]\leq |Q_F|^r=A_3(r)$.

Induct on the length $l$ of an iterated free extension, with all
free ranks at most $r$.  For $l=1$, restriction to $F'F^3$ kills
$H^1(F,k)$, and the other positive cohomology groups vanish.
For $1\to F\to G\xrightarrow{\pi}B\to1$, choose the inductively
constructed subgroup $B_0\leq B$ killing all positive cohomology of
$B$, and set $G_1=\pi^{-1}(B_0)\cap J$.
Its index is at most $a=C(l-1,r)A_3(r)$.
Put $F_1=G_1\cap F$ and $B_1=\pi(G_1)$.
The actual map of Lyndon--Hochschild--Serre spectral sequences is
\[
 H^p\bigl(B,H^q(F,k)\bigr)
       \longrightarrow H^p\bigl(B_1,H^q(F_1,k)\bigr).
                                                               \tag{ECR.10}
\]
Only $q=0,1$ occur.  In the $q=1$ row the coefficient homomorphism
$H^1(F,k)\to H^1(F_1,k)$ is zero, since $F_1\subset F'F^3$.
Thus that entire row map is zero without requiring a trivial base
action on these coefficient modules.  In the $q=0$ row and positive
total degree, restriction factors through $H^p(B_0,k)$ and is zero
by construction.  Consequently the induced map on the associated
graded of each positive total-degree cohomology group is zero.
Its filtration has at most two pieces.  Repeat this same construction
inside $G_1$; all its inherited ranks are at most $1+a(r-1)$.
The composite of the two restriction maps strictly moves the
two-step filtration twice, and is zero on the actual positive
cohomology groups.  Its index is bounded by exactly $C(l,r)$ in
(ECR.2).  This proves the induction, including the displayed bound.

For any subgroup $U\leq K_0$ of index at most $d$, the construction
just proved supplies $U'\leq U$ of index at most $C(8,1+8d)$ whose
restriction kills every positive cohomology group of $U$.
Define
\[
 \kappa_m(K_0)=\bigcap_{\varphi:K_0\to\mathfrak S_m}\ker\varphi.
                                                               \tag{ECR.11}
\]
The actual coset action of $K_0$ on $K_0/U'$ has degree at most
$m=dC(8,1+8d)$; embed a smaller symmetric group into
$\mathfrak S_m$ by fixing the extra letters.  Hence
$\kappa_m(K_0)\subset U'\subset U$.  Restriction from $U$ to
$\kappa_m(K_0)$ is still zero in all positive degrees.
This subgroup is characteristic in $K_0$, and its index is at most
$(m!)^{(m!)^{44}}$, since a homomorphism to $\mathfrak S_m$ is
specified by the images of the $44$ chosen generators.
Starting with $U_0=\ker(q|_{K_0})$, set
\[
 U_{i+1}=\kappa_{m_i}(K_0)\quad(0\leq i<289),\qquad V=U_{289}.
                                                               \tag{ECR.12}
\]
The increasing bounds (ECR.2)--(ECR.3) prove inductively that these
are nested, that $[K_0:U_i]\leq D_i$, and that every positive
cohomology restriction $H^*(U_i,k)\to H^*(U_{i+1},k)$ is zero.
The maps and kernels in (ECR.9)--(ECR.12) are finite constructions.
For example all homomorphisms in (ECR.11) are enumerated from a
finite presentation supplied by the cell model below.

\par\medskip\noindent\textbf{A bounded finite equivariant cell model.}\quad
Write $z_0=0$ and use the sixteen real torus coordinates of the eight
other points.  In their universal affine coordinate space take the
periodic rational hyperplanes
\[
 3\bigl(z_h^{(j)}-z_k^{(j)}\bigr)\in\mathbb Z,
       \qquad h\ne k\in\mathbb F_3^2,\quad j=1,2.
                                                               \tag{ECR.13}
\]
There are $72$ normal directions; a unit cube meets at most seven
levels of each, so this arrangement has a finite polyhedral quotient
on the torus.  The action (NCF.2) permutes it.  Every removed diagonal
is a union of its cells, and all forty-eight fixed grids are vertices.
The coordinate directions are among (ECR.13), so its bounded cells
give the usual genuine polyhedral cell structure before quotienting
by the integral lattice.  Barycentric subdivision makes stabilizers
fix their cells pointwise.  The complement of the diagonal subcomplex
deformation retracts equivariantly onto the full barycentric
subcomplex on cells not contained in a removed diagonal: in each
barycentric simplex delete the coordinates of the removed vertices
and renormalize the remaining coordinates.  At least one remaining
coordinate is positive precisely on this complement.  This formula
is compatible on faces and gives the claimed retraction.

Denote the resulting finite complex by $Y$.  It has dimension at
most sixteen, contains all forty-eight grid fixed points, and
$Y^H=X^H$.  Lift this cell complex and the retraction to every actual
cover specified by $U_i$.  The normality and $H$-invariance of $U_i$
give compatible $H$ actions fixing the chosen basepoint lifts.
These are exactly the actions obtained from the split source
$K_0\rtimes H$.  Write $Y_i$ for the resulting connected covers.
Their cellular cochains, covering pullbacks, and group actions are
finite matrices in degrees zero through sixteen.  A stabilizer of a
lifted cell fixes it pointwise as well, because its projection does
and a deck transformation fixing a point is identity.

\par\medskip\noindent\textbf{The two cyclic Tate comparisons.}\quad
Here is the chain argument, including its finite filtration bound.
For a cellular $C_3=\langle g\rangle$ action on a finite complex
$Z$ of dimension at most sixteen, fixing a chosen vertex $z$, use
$C^*(Z,z;k)$.  Form the complete Tate double complex using the
rank-one complete free resolution with alternating differentials
$1-g$ and $1+g+g^2$.  Its spectral sequence obtained by taking
ordinary cochain cohomology first has
\[
 E_2^{s,t}=\widehat H^s\bigl(C_3,H^t(Z,z;k)\bigr),
                    \qquad 0\leq t\leq16.
                                                               \tag{ECR.14}
\]
Although the horizontal resolution is bi-infinite, every total
degree has only these seventeen vertical terms.  Thus the filtration
and convergence are finite in each total degree.  A map zero on all
ordinary relative cohomology is zero on this $E_2$ page, and hence
strictly moves the induced filtration on Tate hypercohomology.
Seventeen successive such maps have zero composite on the latter.
This uses both alternating differentials of the complete resolution.

Restriction to the fixed subcomplex is an isomorphism on Tate
hypercohomology.  Indeed the relative cellular cochains
$C^*(Z,Z^{C_3};k)$ are sums of free orbit modules in each degree;
their horizontal Tate complexes are exact.  The bounded vertical
range then makes their total Tate complex acyclic.  On the fixed
subcomplex the action is trivial, so both horizontal differentials
are zero over $k$.  Its Tate hypercohomology is the direct sum, in
each total degree, of the ordinary relative cohomology groups with
their corresponding horizontal shifts.  A zero Tate map therefore
gives zero on every ordinary relative cohomology group of the fixed
subcomplex.

Apply this calculation to a chosen subgroup $A\cong C_3$ of $H$.
The spaces $Y_i$ are connected, so the vanishing of positive
cohomology restrictions from (ECR.12) also gives vanishing in every
degree for cohomology relative to the chosen basepoints.
Each block of seventeen covering pullbacks therefore gives a zero
relative cohomology map on the $A$-fixed subcomplexes.  Now apply the
same calculation to $H/A$ acting on those subcomplexes, which still
have dimension at most sixteen.  Seventeen blocks suffice.  The
result after $17^2=289$ covers is
\[
 H^0(Y_0^H,\{y_0\};k)\longrightarrow
                   H^0(Y_{289}^H,\{y_{289}\};k)
                         \quad\hbox{is zero}.                  \tag{ECR.15}
\]
Both fixed sets are finite and discrete, since they project under
actual finite covering maps to $X^H$.  The characteristic function of
any point of $Y_0^H$ other than $y_0$ is a class in the source of
(ECR.15).  Its zero pullback proves that every point of $Y_{289}^H$
maps to $y_0$.  In particular consider the fibre over the chosen base
grid.  In the regular cover it is literally $K_0/U_i$, and $H$ acts
there by its specified conjugation action.  We have proved the exact
finite statement
\[
 \operatorname{im}\left(
 C_{(K_0/V)\rtimes H}(H)\cap K_0/V
                       \longrightarrow K_0/U_0\right)=1.
                                                               \tag{ECR.16}
\]
This is a statement about the displayed quotient homomorphism, not
an inference from the infinite centralizer assertion (NCF.3).

\par\medskip\noindent\textbf{Embedding the quotient in a finite geometric image.}\quad
Because $V$ is characteristic in the normal subgroup $K_0$, each
linear section in (ECR.8) acts on $K_0/V$.  Define the actual finite
homomorphism
\[
 \begin{split}
 \lambda:\mathcal U&\longrightarrow
          Q\times\bigl((K_0/V)\rtimes\operatorname{Aut}(K_0/V)\bigr),\\
 k s(A)&\longmapsto
       \bigl(q(k s(A)),\ (kV,\operatorname{Ad}_{s(A)}|_{K_0/V})\bigr).
 \end{split}                                                    \tag{ECR.17}
\]
Its image $R$ has order at most $D [K_0:V]^{45}$, since the
automorphism group of the $44$-generated finite group $K_0/V$ has
order at most $[K_0:V]^{44}$.  Its kernel intersects $K_0$ exactly
in $V$ and is contained in $\ker(q|_{\mathcal U})$.

Induce (ECR.17) through the $I$ actual cosets of the normal subgroup
$\mathcal U$ in $\mathcal M$, retaining the source-word weights and
the coset permutation.  Let $\eta$ be this homomorphism with its
actual finite image $E$.  Its permutation image has order $I$, so
\[
 |E|\leq I|R|^I\leq I\bigl(D D_{289}^{45}\bigr)^I.
                                                               \tag{ECR.18}
\]
The kernel of induction is the intersection of the conjugates of
$\ker\lambda$.  Its intersection with $K_0$ is precisely $V$,
because $V$ is preserved by all of $\mathcal M$.  Thus
$\eta(K_0)=K_0/V$ by the specified identity map.  At the identity
coset induction recovers $\lambda$, so $\ker\eta\subset\ker q$;
this defines the literal projection $p$ in (ECR.4).
The permutation part retains the faithful labels of $H$.

\par\medskip\noindent\textbf{The exact normalizer projection.}\quad
Take $n\in N_E(\eta H)$ and choose an actual source representative
$g\in\mathcal M$.  Its label permutation is affine,
$h\mapsto\nu h+t$, because it normalizes the regular translation
permutations.  Its image under $p$ belongs to $N_Q(qH)$.
The retained lower comparison of A.4.24--A.4.25 therefore gives
$\nu=A\bmod3$, where $A$ is the integral aligned compact period
of $g$.  In particular $\nu$ has determinant one.
The actual affine torus map
\[
                  \ell:z\longmapsto Az+t/3
                                                               \tag{ECR.19}
\]
has the same period and labels as $g$ and normalizes the original
named $H$.  Both $n$ and $\eta(\ell)$ induce $h\mapsto\nu h$ on
$\eta H$: its labels are faithful.  Consequently
$k=g\ell^{-1}\in K_0$ has $\eta(k)\in C_E(\eta H)$.
Since $\eta|_{K_0}$ has kernel $V$, this says exactly that $kV$
centralizes $H$ in the source of (ECR.16).  That formula gives
$q(k)=1$, whence
\[
                         p(n)=q(\ell).
                                                               \tag{ECR.20}
\]
This proves the first inclusion in (ECR.6), and the opposite inclusion
is given by the actual images of ordinary normalizer elements.
For $n$ in the centralizer, its affine linear label part is identity.
Then $A\equiv I\pmod3$, so the same $\ell$ centralizes every grid
translation.  This proves the centralizer equality as well.
Finally the ordinary zero-label normalizer comparison (MFP.4)
identifies $d\operatorname{pr}_G$ with the aligned ordinary
forgetting map; (ECR.20) transfers that equality through the specified
projection, proving the final assertion.
\end{proof}

\par\medskip\noindent\textbf{A literal period coordinate for subsequent transport.}\quad
The quotient above also gives a convenient exact way to identify the
configuration subgroup in a finite transport domain.  Restrict
$\eta\circ s$ to $\Gamma(3)$, induce this actual homomorphism through
its twenty-four cosets in $\operatorname{SL}_2(\mathbb Z)$, and call
the resulting finite period map $\Omega$.  It retains reduction
modulo three and has order at most $24|E|^{24}$.  Form
\[
 \eta^*=(\eta,\Omega\circ\operatorname{per}),\qquad
                         E^*=\operatorname{im}\eta^*.
                                                               \tag{ECR.21}
\]
Within $E^*$, the elements having a translation label and identity
$\Omega$ coordinate form exactly
\[
                      (K_0/V)\rtimes H.
                                                               \tag{ECR.22}
\]
To verify the reverse inclusion explicitly, take a source
representative $g$, remove its actual translation $h$ on the right,
and write $gh^{-1}=k s(A)$ using (ECR.8).  The identity $\Omega$
coordinate gives $\eta(s(A))=1$ in the identity induced coordinate.
Thus $\eta^*(g)=\eta^*(k h)$.  The forward inclusion is immediate
from the actual periods and labels of $K_0$ and $H$.
This supplies an actual finite subgroup and its source map for a
subsequent competing-complement calculation.  The theorem above
recognizes the canonical $H$; it does not by itself assert fusion of
every competing finite complement under a transported action.
